\section{Materials and Methods}
\label{sec:materials_and_methods}

\begin{figure*}[htbp]
\centering
\includegraphics[width=\textwidth]{method.png}
\caption{End-to-End Framework of the Proposed Multi-Domain Vision Transformer Feature Fusion and Classical Classifier Optimization for Intersectional Race and Gender Classification.}
\label{fig:framework}
\end{figure*}

The proposed framework integrates multi-domain visual representations and machine learning model optimization for intersectional race and gender classification, as illustrated in Figure~\ref{fig:framework}. The methodology begins with partitioning $224 \times 224$ pixel resolution facial images from the DemogPairs dataset into training and test partitions via a stratified split. Three frozen pre-trained ViT-Base models extract 768-dimensional latent representations from the [CLS] token offline. The three vectors are combined through concatenation fusion $\mathbf{z}_{\text{tri}} = \mathbf{f}_{\text{face}} \oplus \mathbf{f}_{\text{emotion}} \oplus \mathbf{f}_{\text{age}}$, where $\mathbf{z}_{\text{tri}}$ denotes the 2,304-dimensional tri-domain vector, $\oplus$ is the concatenation operator, and $\mathbf{f}_{\text{face}}$, $\mathbf{f}_{\text{emotion}}$, and $\mathbf{f}_{\text{age}}$ represent facial biometric, expression, and age features, respectively. The representation vectors are processed within a modular Scaler and PCA pipeline before being classified into six intersectional classes by RF, GNB, LR, and SVM. Hyperparameter tuning is conducted via 5-Fold Stratified Cross-Validation with an isolated GridSearchCV scheme to prevent information leakage prior to test data evaluation.
